\section{System's Perspective}\label{systems-perspective}

This section will give description and illustrations of the design and
architecture of our Mini-twit system, dependencies and the current state
of out system. \\

The application is based on Chirp - a c\# application made in the course Analysis, Design and Software Architecture at ITU. 
We chose to base the project on Chirp as all of knew the codebase.


\subsection{Design and architecture of our ITU-MiniTwit
systems.}\label{design-and-architecture-of-our-itu-minitwit-systems.}

The architecture of MiniTwit follows, as Chirp, the onion architecture. 
The deployment architecture of the MiniTwit system can be seen in figure UML\_DIAGRAM.
The architecture strives to achieve both high-availability and the possibility for scaling when needed. 
A user can reach MiniTwit with the domain \texttt{Devbobs.tech} on port 433 or 80, which redirects to 443.
The request will reach the Leader node and will then get loadbalanced between the four instances of MiniTwit running.
MiniTwit is also available on port 8080 with http. The request will in this case be loadbalanced by the docker ingress routing mesh.
The architecture have a single point of failure, as will be explained in \ref{availability}.




\subsection{Dependencies}\label{dependencies}

package and deliver services = docker containers (has all dependencies).
Nuget dependencies. Github action dependencies (Hadolint, Csharpier, ).
Docker.

\subsection{Current state of system}\label{current-state-of-system}
